% Independent extension of ISR. All original labels are retained.
\section{Positivity eliminates the deep defining-prime strata}
Retain the actual positive ES hypotheses of ISR1 and the classification
proved in ISR2--3. Permuting two denominator labels below is explicitly
only for the inequality proof; all original ordered matrices are returned
by the same permutation of their rows and branches. No two root or factor
sign labels are identified.

\subsection{The two unit denominators differ by less than $p^2$}
In the one-divisible case write $z=pZ$ and order the two unit denominators
as $x<y$. ISR2 gives $Z\equiv1/4\pmod p$. Since $p\equiv1\pmod4$,
the least positive residue of $1/4$ is $(3p+1)/4$. Positivity and the
strict inequality $1/x<4/p$ therefore imply
\begin{equation}
 Z\ge Z_0=(3p+1)/4,\qquad x\ge x_0=(p+3)/4,\qquad
 y=\frac{pZx}{Z(4x-p)-x}.                                \tag{PES1}
\end{equation}
The denominator at $(x_0,Z_0)$ is $2p>0$. On the entire rectangle
$x\ge x_0,Z\ge Z_0$, its derivative in $x$ is $4Z-1>0$ and its
derivative in $Z$ is $4x-p>0$, so the denominator remains positive.
For its exact gap function $g(x,Z)=y-x$, differentiation gives
\begin{equation}
 \partial_x g=-\frac{p^2Z^2}{(Z(4x-p)-x)^2}-1<0,\qquad
 \partial_Z g=-\frac{p x^2}{(Z(4x-p)-x)^2}<0.             \tag{PES2}
\end{equation}
Consequently every actual one-divisible witness satisfies
\begin{equation}
 y\le\frac{(p+3)(3p+1)}{32},\qquad
 0<y-x\le\frac{(p+3)(3p-7)}{32}<p^2,
 \quad v_p(y-x)\le1.                                  \tag{PES3}
\end{equation}
The last strict inequality is equivalent to
$29p^2-2p+21>0$. These bounds do not assume $p\nmid S$.
The real upper gap bound is attained by the actual witness
$(p;x,y,z)=(13;4,20,130)$, for which $y-x=16$.
This is sharpness of the displayed real inequality; it does not
assert that valuation one occurs among actual witnesses.

For later finite searches, ordering $x\le y$ also gives
$2/x\ge4/p-1/(pZ)$, hence
\begin{equation}
 x\le\frac{2pZ}{4Z-1}
   =\frac p2+\frac{p}{2(4Z-1)}
   \le\frac p2+\frac16,
 \quad x\le(p-1)/2.                                    \tag{PES4}
\end{equation}
All these inequalities are consequences of the original reciprocal
equation and positivity. They exclude every formal one-divisible
Smith stratum with $n\ge2$ from actual positive witnesses.

\subsection{Every two-divisible witness has separated scaled residues}
In the two-divisible case order $1\le Y<Z$ after retaining the
corresponding permutation of $(y,z)=(pY,pZ)$. Put $r=4x-p>0$.
The exact reciprocal equation and ISR3 give
\begin{equation}
 \frac r x=\frac1Y+\frac1Z,\quad
 Y+Z\equiv4YZ\pmod p,
 \quad Y\le\frac{2x}{4x-p}\le\frac{p+3}{6}<p.           \tag{PES5}
\end{equation}
The second inequality follows because $2x/(4x-p)$ has negative
derivative $-2p/(4x-p)^2$ and $x\ge(p+3)/4$.
The root distinctness proof in ISR1 implies $Y,Z\ne1$ as integers.
There can be no repeated residue among $1,Y,Z$:

If $Y\equiv1\pmod p$, its positive bound $Y<p$ forces $Y=1$,
contradicting the preceding fact. If $Z\equiv1\pmod p$, the reduced
equation forces $3Y\equiv1$. Since $p\equiv1\pmod3$, its least
positive solution is $(2p+1)/3$, larger than $(p+3)/6$.
Finally, if $Y\equiv Z\pmod p$, the same equation and their unit
property imply $2Y\equiv1$. Its least positive solution is
$(p+1)/2$, again larger than $(p+3)/6$. Therefore
\begin{equation}
 v_p(p-pY)=v_p(p-pZ)=v_p(pY-pZ)=1.                       \tag{PES6}
\end{equation}
This proves that the closer-pair parameter in ISR17 is always $n=1$
for an actual positive hard-prime witness. In particular $p\nmid S$
automatically, and the entire invariant-factor list and conductor data
are universal on this actual class:
\begin{equation}
 \begin{aligned}
 &\operatorname{Smith}(O/E)=\operatorname{Smith}(E/I)
       =(0,0,0,1,1,2,2,3),\qquad L=9,\\
 &(b_p,b_x,b_y,b_z)=(2,0,2,2),\qquad
 (m_p,m_x,m_y,m_z)=(1,0,1,1),\\
 &(\epsilon_p,\epsilon_x,\epsilon_y,\epsilon_z)=(0,0,0,0),
 \qquad (v_p I_i)=(3,0,3,3).
 \end{aligned}                                         \tag{PES7}
\end{equation}
Thus every normalized quadratic branch is unramified or split; the
nonreduced length-six clustered integral fibre remains present.

\subsection{The earlier coefficient ideal always enters the actual ideal}
For any actual positive witness with $p\nmid S$, the criterion PS16
now holds without an additional gap hypothesis. In the one-divisible
case $v_p(C_0)=1$ and PES3 gives $n\in\{0,1\}$, so the largest
branch exponent $3m_i+\epsilon_i$ is one. In the two-divisible case
$v_p(C_0)=2$ and PES7 gives largest branch exponent three. Hence
\begin{equation}
 C_0^2 E\subset I
 \quad\text{for every actual positive witness with }p\nmid S.
                                                               \tag{PES8}
\end{equation}
This is the exact specialization of the coefficient conductor
from PS14--16, not an equality of the two ideals. Both original
complex receiving maps of ISR32--34 apply unchanged to this proven
inclusion. For $p\mid S$ the original coefficient columns are the
fractional lattice of ISR18--23, so PES8 is not incorrectly stated
as an inclusion of an undefined original integral order.

\subsection{Exact parameters for the remaining two bounded questions}
The proofs above do not establish that the remaining one-divisible
case $n=1$ is realizable or impossible, and do not prove that $p\mid S$
is impossible. The following exact derivation constrains both questions.
Write $x=g a$, $y=g b$ with $a<b$, $\gcd(a,b)=1$ and $p\nmid gab$.
Set $d=4gab-p(a+b)>0$. From the ES equation,
$Z=gab/d$. The integer $d$ is coprime to both $a$ and $b$ because
its reductions modulo them are $-pb$ and $-pa$. Thus $d\mid g$.
Write $g=hd$. Rearranging the definition of $d$ gives
$d(4hab-1)=p(a+b)$. Also $p\nmid d$, since
$d\equiv4gab\not\equiv0\pmod p$. Consequently there are positive
integers $k,r$ with the complete identities
\begin{equation}
 \begin{aligned}
 &x=hda,\quad y=hdb,\quad Z=hab,\quad
 k=(4hab-1)/p,\quad a+b=dk,\\
 &r=4x-p=4had-p,\quad rk=4ha^2+1,\quad
 b=dk-a,\quad p=4had-r,\quad r\equiv k\equiv3\pmod4.
 \end{aligned}                                         \tag{PES9}
\end{equation}
For $rk$, substitute $b=dk-a$ into $pk=4hab-1$ and cancel to obtain
$(4had-p)k=4ha^2+1$. The two residue classes modulo four follow
from $p\equiv1\pmod4$ and the product $rk\equiv1\pmod4$.
No parameter has been replaced by a unit representative.
The two still possible residue events have exact tests
\begin{equation}
 \begin{aligned}
 x\equiv y\pmod p
   &\Longleftrightarrow p\mid(dk-2a),
       &&\Longrightarrow p\mid(4ha^2-1),\\
 p\mid S
   &\Longleftrightarrow p\mid k,
       &&\Longrightarrow p\mid(4ha^2+1).
 \end{aligned}                                         \tag{PES10}
\end{equation}
The equivalences use $p\nmid hd$ and $S\equiv x+y\equiv hd^2k$.
For the first implication, $a\equiv b$ and $4hab-1=pk$ give
$4ha^2-1=0\pmod p$. The second follows from $rk=4ha^2+1$.
These are exact necessary conditions and a finite-search parametrization;
neither event is excluded merely because a bounded search finds none.

\subsection{Every inclusion entry and the least coefficient power}
Let $c=v_p(C_0)$ and keep the complete original ideal basis $G_I$ of
ISR12. From $VQ=\operatorname{diag}(f'(r_i))$, the full eight-column
factorization of the scalar ideal inclusion is
\begin{equation}
 Z_j=G_I^{-1}C_0^j I_8
   =\operatorname{diag}\left(
       \operatorname{diag}\left(\frac{C_0^j}{p^{m_i}D_i}\right)V,
       \operatorname{diag}\left(\frac{C_0^j}{D_i}\right)V\right),
 \quad G_IZ_j=C_0^jI_8.                                \tag{PES11}
\end{equation}
Thus its complete even block entry $(i,n)$ is
$C_0^j r_i^n/(p^{m_i}D_i)$ and its complete odd block entry is
$C_0^j r_i^n/D_i$, for $0\le i,n\le3$. Because each original root
is integral and $m_i\ge0$, all entries are integral if
$jc\ge\max_i(m_i+b_i)$. Conversely the even entries with $n=0$
force exactly these four inequalities. Therefore the least power
and both unchanged receiving identities are
\begin{equation}
 \min\{j\ge0:C_0^jE\subset I\}
  =\begin{cases}1&\text{one divisible denominator},\\
                 2&\text{two divisible denominators},\end{cases}
 \qquad (F_U G_I)Z_j=C_0^jF_U,\quad
        (F_W G_I)Z_j=C_0^jF_W.                          \tag{PES12}
\end{equation}
Here $c=1$ and the largest $m_i+b_i$ equals one in the first case;
$c=2$ and the largest equals three in the second. The matrices
$F_U,F_W$ retain all original moments, centres and Gamma masses
from ISR32--34. This proves every assertion of PB6--8 through the
entire original coefficient matrices, including necessity of the
least exponents, not just an index comparison.

\paragraph{Provenance.}
The ES relation, actual-witness marking, and finite signed algebra
retain the human-source citations and original programme definitions
of ISR1 and PS1. PES1--12 are the independent positivity and parameter
derivations written above. No new external literature is claimed read,
and no Frobenius or RH conclusion is assigned to this finite calculation.
